%% notes-tex/028-gene-essentiality/tables/round_w67.tex
%% GENERATED by experiments/028-gene-essentiality/scripts/notes_tex_tables.py
%% SOURCE: results/w67_readout.json
\begin{table}[H]
  \centering
  \scriptsize
  \caption[Round 51, the hash seed fixed on the old one-hot code]{Round 51, the hash seed fixed on the old one-hot code. The round-9 pooled cell (pathway multi-hot and ProtT5, the plain pooled head) and the round-33 layer-free cell, each run twice on seed 2700 with the card's kernels pinned and Python's hash seed fixed, on the code that rebuilt the pathway one-hot from a set at every launch. Batch 32, class weight, 200 epochs; the read is per pair. Bold marks the highest value in a column. $\uparrow$ ($\downarrow$) marks an arm whose held-out AUROC is above (below) the first arm's on every one of the same three seeds.}
  \label{tab:w67}
  \begin{tabular}{@{}p{52mm}rrlllr@{}}
    \hline
    arm & mean & sd & seeds & pooled [95\% CI] & selected epochs & val step \\
    \hline
    round-9 pooled cell, kernels pinned, PYTHONHASHSEED=0, old one-hot code, first run & 0.860 & 0.000 & 0.860 & 0.860 [0.79, 0.93] & 31 & 0.0036 \\
    the same, second run & 0.853 & 0.000 & 0.853 & 0.853 [0.78, 0.92] & 54 & 0.0039 \\
    round-33 layer-free cell, kernels pinned, PYTHONHASHSEED=0, old one-hot code, first run & \textbf{0.913} & 0.000 & 0.913 & \textbf{0.913} [0.87, 0.95] & 39 & 0.0039 \\
    the same, second run & \textbf{0.913} & 0.000 & 0.913 & \textbf{0.913} [0.87, 0.95] & 39 & 0.0039 \\
    \hline
  \end{tabular}
\end{table}
